% [EN] Build: cd docs/reports/reconstruction && latexmk -pdf nn_retrain_qmdwindow_B115T3deg.tex / [CN] 编译: cd docs/reports/reconstruction && latexmk -pdf nn_retrain_qmdwindow_B115T3deg.tex
\documentclass[11pt]{beamer}
\usetheme{Madrid}
\usepackage{booktabs}
\usepackage{amsmath}

\title{QMD-Window Retraining Plan and Results\\B115T, Target 3deg}
\author{SMSimulator Reconstruction}
\date{\today}

\begin{document}

\begin{frame}
\titlepage
\end{frame}

\begin{frame}{Retraction / Update --- 2026-04-20}
\textbf{Contamination found in original qmd-window training}
\begin{itemize}
\item The original \texttt{formal\_B115T3deg\_qmdwindow/20260227\_223007/}
training consumed Geant4 output produced while the
\texttt{DeutDetectorConstruction::fSetTarget} default was silently
\texttt{true}. Tree-gun input (already post-IMQMD secondaries) was
therefore scattered a second time by the Sn target.
\item Audit: \texttt{docs/audit/2026-04-20-nn\_target\_contamination\_audit.md}.
\item The polluted formal run and all 7 \texttt{optimization\_runs/*}
sweeps were deleted; the Results / Sweep / Best-Candidate /
Final-Accuracy tables in this deck therefore no longer reflect
production state.
\end{itemize}

\textbf{Clean retrain:} \texttt{data/nn\_target\_momentum/formal\_B115T3deg/20260420\_184345/}
(same narrow sampling: \(R=100\) MeV/c, \(p_z\in[560,680]\) MeV/c)
\begin{table}
\centering
\small
\begin{tabular}{lccc}
\toprule
Val RMSE (MeV/c) & Contaminated & Clean retrain & Change \\
\midrule
\(p_x\) & 13.60 & \textbf{6.70} & \(-51\%\) \\
\(p_y\) & 12.44 & \textbf{2.40} & \(-81\%\) \\
\(p_z\) & 5.89 & 5.80 & \(-2\%\) \\
\(|\vec p|\) & 5.88 & 6.08 & \(+3\%\) (\(p_z\) dominated) \\
\bottomrule
\end{tabular}
\end{table}
Large transverse gain matches the physics of the bug: Sn multiple
scattering perturbs \(p_x,p_y\) more than \(p_z\).
\end{frame}

\begin{frame}{Motivation}
\begin{itemize}
\item Task: reconstruct target proton momentum \((p_x,p_y,p_z)\) from PDC two-point track.
\item Existing formal model (wide sampling) is usable, but dataset efficiency is low.
\item Baseline campaign: \texttt{20260227\_205732}.
\end{itemize}

\textbf{Baseline key numbers}
\begin{itemize}
\item Sampling: \(p_\perp \le 150\) MeV/c, \(p_z \in [500,700]\) MeV/c.
\item Effective keep fraction (after PDC track reconstruction cuts): \(\sim 24\%\).
\item Momentum-norm RMSE: val \(9.19\) MeV/c, test \(9.84\) MeV/c.
\end{itemize}
\end{frame}

\begin{frame}{Retraining Campaign Definition}
\textbf{Geometry and setup}
\begin{itemize}
\item Geometry macro:
\texttt{configs/simulation/DbeamTest/detailMag1to1.2T/geometry\_B115T\_pdcOptimized\_20260227.mac}
\item Target transform in run macro:
\((-1.2488,\,0.0009,\,-106.9458)\) cm, angle \(3.0^\circ\).
\end{itemize}

\textbf{New sampling window}
\begin{itemize}
\item \(p_\perp \le 100\) MeV/c (disk in \(p_x,p_y\)).
\item \(p_z \in [560,680]\) MeV/c.
\end{itemize}

\textbf{Split size and seeds}
\begin{itemize}
\item train/val/test = \(200000/30000/30000\).
\item Seeds: 20260301 / 20260302 / 20260303.
\end{itemize}
\end{frame}

\begin{frame}{Pipeline}
\begin{enumerate}
\item Generate tree-input ROOT with proton truth sampled in new momentum window.
\item Run Geant4 simulation for train, val, test independently.
\item Build dataset CSV/ROOT with PDC track reconstruction.
\item Train MLP: \(6 \rightarrow 128 \rightarrow 128 \rightarrow 64 \rightarrow 3\).
\item Evaluate val/test with independent hold-out files.
\end{enumerate}

\textbf{Tools}
\begin{itemize}
\item \texttt{run\_formal\_training.sh}
\item \texttt{generate\_tree\_input\_disk\_pz.C}
\item \texttt{build\_dataset.C}, \texttt{train\_mlp.py}, \texttt{infer\_mlp.py}, \texttt{evaluate\_mlp.py}
\end{itemize}
\end{frame}

\begin{frame}{MLP Architecture}
\textbf{Network: \(6 \rightarrow 128 \rightarrow 128 \rightarrow 64 \rightarrow 3\)}
\begin{itemize}
  \item Input (6): \texttt{pdc1\_x, pdc1\_y, pdc1\_z, pdc2\_x, pdc2\_y, pdc2\_z}
  \item Three hidden layers with \textbf{ReLU} activations: 128, 128, 64 units.
  \item Output (3): \((p_x, p_y, p_z)_\text{truth}\) in MeV/c (no output activation).
\end{itemize}

\textbf{Evaluation metrics (post-training)}
\begin{itemize}
  \item Per-component RMSE, MAE, bias, and \(p_{95}\) for \(p_x, p_y, p_z\).
  \item Momentum-norm RMSE: \(\sqrt{\langle(|\hat{\vec{p}}| - |\vec{p}|)^2\rangle}\).
\end{itemize}
\end{frame}

\begin{frame}{Loss Function and Optimiser}
\textbf{Training loss: MSE over all three components}
\[
  \mathcal{L} = \frac{1}{3N}\sum_{i=1}^{N}
  \left[(\hat{p}_x^{(i)}-p_x^{(i)})^2
       +(\hat{p}_y^{(i)}-p_y^{(i)})^2
       +(\hat{p}_z^{(i)}-p_z^{(i)})^2\right]
\]
Implemented as \texttt{nn.MSELoss()} in PyTorch.

\textbf{Optimiser and hyperparameters}
\begin{itemize}
  \item Optimiser: \textbf{Adam}, \(\text{lr}=10^{-3}\), weight decay \(10^{-5}\).
  \item Batch size: 1024; Max epochs: 120.
  \item \textbf{Early stopping}: patience = 12 epochs on validation loss;
        best-checkpoint model is restored at the end.
\end{itemize}
\end{frame}

\begin{frame}{Data Preprocessing and Split}
\textbf{Feature normalisation (inputs only)}
\begin{itemize}
  \item \(z\)-score standardisation using training-set mean and std:
        \(\tilde{x} = (x - \mu_\text{train})/\sigma_\text{train}\).
  \item Labels \((p_x, p_y, p_z)\) are \emph{not} normalised.
  \item Mean/std saved in \texttt{model\_meta.json} for inference.
\end{itemize}

\textbf{Dataset split (external files, fixed seeds)}
\begin{itemize}
  \item train / val / test = 200\,000 / 30\,000 / 30\,000 events.
  \item Seeds: 20260301 / 20260302 / 20260303.
  \item Val loss used for model selection; test set held out until final evaluation.
\end{itemize}
\end{frame}

\begin{frame}{Acceptance and Quality Criteria}
\textbf{Primary objective}
\begin{itemize}
\item Improve effective kept samples by reducing out-of-acceptance phase space.
\end{itemize}

\textbf{Evaluation criteria}
\begin{itemize}
\item Keep fraction should improve from baseline \(\sim24\%\).
\item Test \(|\vec{p}|\) RMSE should remain stable or improve.
\item Component-wise RMSE (\(p_x,p_y,p_z\)) tracked for regression.
\end{itemize}
\end{frame}

\begin{frame}{Results \textcolor{red}{[CONTAMINATED --- see retraction frame]}}
\begin{table}
\centering
\begin{tabular}{lcc}
\toprule
Metric & Baseline (Wide) & QMD-Window (New) \\
\midrule
Train kept fraction & 23.61\% & 44.48\% \\
Val kept fraction & 23.14\% & 44.16\% \\
Test kept fraction & 23.44\% & 45.02\% \\
Val \(|\vec{p}|\) RMSE (MeV/c) & 9.19 & 6.01 \\
Test \(|\vec{p}|\) RMSE (MeV/c) & 9.84 & 5.88 \\
\bottomrule
\end{tabular}
\end{table}
\footnotesize
\textcolor{red}{RMSE numbers from the polluted model. Kept-fraction numbers
(geometry acceptance only) are still valid. Clean retrain: val \(|\vec p|\) RMSE
\(= 6.08\); see retraction frame for component-wise detail.}
\end{frame}

\begin{frame}{Post-Baseline Optimisation Objective}
\textbf{Why continue after qmd-window?}
\begin{itemize}
\item QMD-window already improved usable sample rate and hold-out \(|\vec{p}|\) RMSE.
\item The next target is narrower: further reduce proton \(p_x/p_y\) reconstruction error on hold-out data.
\item Project ranking metric for follow-up runs:
\[
\mathrm{pxpy\_rmse\_mean}=\frac{\mathrm{RMSE}(p_x)+\mathrm{RMSE}(p_y)}{2}.
\]
\end{itemize}

\textbf{Search axes actually explored}
\begin{itemize}
\item Manual loss weighting to emphasize \(p_x,p_y\).
\item Target \(z\)-score normalization.
\item Wider MLP hidden dimensions.
\end{itemize}

\textbf{Important comparison rule}
\begin{itemize}
\item After target normalization, raw \texttt{best\_val\_loss} changes scale.
\item Therefore cross-run ranking is based on raw hold-out RMSE metrics, not on loss value alone.
\end{itemize}
\end{frame}

\begin{frame}{Optimisation Progress: Loss and Normalisation Sweep \textcolor{red}{[CONTAMINATED]}}
\small
\textcolor{red}{All 7 sweep runs consumed the polluted dataset; the sweep was
deleted and is not re-run --- a clean re-sweep may pick a different winner.}
\begin{table}
\centering
\begin{tabular}{lccc}
\toprule
Run & Val \(pxpy\) RMSE & Val \(p_z\) RMSE & Test \(|\vec p|\) RMSE \\
\midrule
Baseline qmd-window & 13.0186 & 5.8863 & 5.8798 \\
\(2,2,1\) & 12.9918 & 5.9981 & 5.9780 \\
\(1.5,1.5,1\) & 12.9977 & 6.1136 & 6.1094 \\
\(3,3,1\) & 12.9235 & 5.9605 & 5.9649 \\
\(3,3,1\) + zscore & 12.8254 & 5.4072 & 5.3718 \\
\bottomrule
\end{tabular}
\end{table}

\textbf{Reading of the sweep}
\begin{itemize}
\item Weighted loss alone helps, but gains remain modest.
\item Target normalization produces the first clearly stronger jump in both \(pxpy\) and \(|\vec p|\) metrics.
\end{itemize}
\end{frame}

\begin{frame}{Current Best Candidate and Interpretation \textcolor{red}{[CONTAMINATED]}}
\textcolor{red}{Both rows of the table below are from the polluted dataset
and are retained only to document the pre-audit state.}
\textbf{Reference vs current best training candidate}
\begin{table}
\centering
\begin{tabular}{lcc}
\toprule
Metric & Baseline qmd-window & Best current candidate \\
\midrule
Configuration & 128-128-64, MSE & 256-256-128, \(4,3,1\)+zscore \\
Val \(p_x\) RMSE & 13.5973 & 13.3366 \\
Val \(p_y\) RMSE & 12.4398 & 12.2689 \\
Val \(p_z\) RMSE & 5.8863 & 5.3517 \\
Val \(pxpy\) RMSE mean & 13.0186 & 12.8027 \\
Test \(|\vec p|\) RMSE & 5.8798 & 5.2488 \\
\bottomrule
\end{tabular}
\end{table}

\textbf{Observed gain}
\begin{itemize}
\item \(pxpy\) RMSE mean improves by about \(1.66\%\).
\item Test \(|\vec p|\) RMSE improves by about \(10.73\%\).
\item Width scaling still helps a little, but gains are already small after normalization.
\item Next likely leverage is better objective design or data-domain refinement, not unlimited width growth.
\end{itemize}
\end{frame}

\begin{frame}{Data Flow (Target-Corrected Full-Event Campaign)}
\textbf{Execution chain}
\begin{enumerate}
\item Reuse existing QMD-derived Geant4 input trees under:
\texttt{data/simulation/g4input/y\_pol/phi\_random/...} and
\texttt{data/simulation/g4input/z\_pol/b\_discrete/...}
\item Run campaign script:
\texttt{scripts/analysis/run\_sn124\_nn\_reco\_pipeline.sh}
\item Script generates per-file Geant4 macros under:
\texttt{data/simulation/g4output/eval\_20260227\_235751/\_macros/*.mac}
\item Each generated macro executes:
\texttt{configs/simulation/DbeamTest/detailMag1to1.2T/geometry\_B115T\_pdcOptimized\_20260227\_target3deg.mac}
\item Geant4 output ROOT and NN reconstruction ROOT are written and then evaluated.
\end{enumerate}
\end{frame}

\begin{frame}{Where Artifacts Are Written}
\textbf{Geometry and model}
\begin{itemize}
\item Geometry macro (target-fixed 3deg, \texttt{SetTarget false}):\\
\texttt{configs/simulation/DbeamTest/detailMag1to1.2T/geometry\_B115T\_pdcOptimized\_20260227\_target3deg.mac}
\item \textbf{Clean production model (2026-04-20 retrain):}\\
\texttt{data/nn\_target\_momentum/formal\_B115T3deg/20260420\_184345/model/model.pt}\\
\textcolor{red}{Previous path
\texttt{formal\_B115T3deg\_qmdwindow/20260227\_223007/model/} was deleted
due to target contamination.}
\end{itemize}

\textbf{Campaign outputs}
\begin{itemize}
\item Full Geant4 outputs:\\
\texttt{data/simulation/g4output/eval\_20260227\_235751/}
\item Full reconstruction outputs (polluted model, kept as baseline):\\
\texttt{data/reconstruction/eval\_20260227\_235751/}
\item Full reconstruction outputs (clean model, 2026-04-22 rerun):\\
\texttt{data/reconstruction/eval\_20260422\_cleanmodel/}
\item Smoke-test outputs:\\
\texttt{data/simulation/g4output/eval\_target3deg\_smoke\_20260228\_\{ypol,zpol\}}\\
\texttt{data/reconstruction/eval\_target3deg\_smoke\_20260228\_\{ypol,zpol\}}
\end{itemize}
\end{frame}

\begin{frame}{Final Accuracy (Clean Model, All Files, All Events)}
\small
\textcolor{gray}{Re-evaluated 2026-04-22 with the clean model over the same
\texttt{eval\_20260227\_235751/} Geant4 outputs; reconstruction step re-run
into \texttt{eval\_20260422\_cleanmodel/}.}
\begin{table}
\centering
\begin{tabular}{lccc}
\toprule
Metric & Overall & Y-pol & Z-pol \\
\midrule
Files & 156 & 16 & 140 \\
Events total & 5,413,747 & 5,178,328 & 235,419 \\
Matched events & 3,632,185 & 3,463,201 & 168,984 \\
Reco efficiency (\%) & 67.092 & 66.879 & 71.780 \\
\(p_x\) RMSE (MeV/c) & \textbf{18.328} & \textbf{17.403} & \textbf{31.828} \\
\(p_y\) RMSE (MeV/c) & 11.170 & 11.141 & 11.746 \\
\(p_z\) RMSE (MeV/c) & 16.072 & 15.073 & 29.937 \\
\(|\vec{p}|\) RMSE (MeV/c) & 13.021 & 12.296 & 23.356 \\
\(p_x\) Bias (MeV/c)  &  -2.016 &  -1.425 & -14.125 \\
\(p_y\) Bias (MeV/c)  &   0.091 &   0.050 &   0.932 \\
\(p_z\) Bias (MeV/c)  &  -6.336 &  -6.928 &   5.786 \\
\bottomrule
\end{tabular}
\end{table}
\end{frame}

\begin{frame}{Clean vs Polluted (RMSE change over full eval)}
\small
\begin{table}
\centering
\begin{tabular}{llrrr}
\toprule
Set & Component & Polluted & Clean & $\Delta$ \\
\midrule
Overall & $p_x$  & 38.73 & \textbf{18.33} & \textbf{-53\%} \\
Overall & $p_y$  & 11.29 & 11.17 & -1\% \\
Overall & $p_z$  & 12.00 & 16.07 & +34\% \\
Overall & $|\vec p|$ &  8.43 & 13.02 & +54\% \\
Y-pol   & $p_x$  & 37.88 & \textbf{17.40} & \textbf{-54\%} \\
Y-pol   & $p_z$  & 11.11 & 15.07 & +36\% \\
Y-pol   & $|\vec p|$ &  7.90 & 12.30 & +56\% \\
Z-pol   & $p_x$  & 53.15 & \textbf{31.83} & \textbf{-40\%} \\
Z-pol   & $p_z$  & 23.81 & 29.94 & +26\% \\
Z-pol   & $|\vec p|$ & 15.73 & 23.36 & +49\% \\
\bottomrule
\end{tabular}
\end{table}
\end{frame}

\begin{frame}{Reconstruction Effect}
\textbf{What changed after fixing target macro}
\begin{itemize}
\item The campaign now runs with explicit target position and target angle in geometry macro.
\item Full-event processing is enforced in production (\texttt{beamOn=0}, all selected files).
\item Reconstruction yields high usable statistics:
\(3{,}632{,}185\) matched proton events from \(5{,}413{,}747\) truth proton events
(\(67.09\%\) matching efficiency --- geometry-only, unaffected by the model swap).
\end{itemize}

\textbf{Clean-model performance (2026-04-22 re-eval)}
\begin{itemize}
\item \(p_x\) RMSE drops \(40\%\)--\(54\%\) across all three sets, matching
the validation-set \(p_x\) \(-51\%\) from the clean retrain.
\item \(p_y\) RMSE flat (\(\sim 11\) MeV/c), bias \(<1\) MeV/c --- NN learns \(p_y\) best.
\item \(p_z\) and \(|\vec p|\) RMSE rise \(26\%\)--\(56\%\); driven by events whose
truth \(p_z\) falls outside the training window \([560, 680]\) MeV/c, where
the polluted model was accidentally better.
\item \(p_z\) bias: overall \(-6.3\), y-pol \(-6.9\), z-pol \(+5.8\) MeV/c. Needs
per-truth-bin residual breakdown.
\end{itemize}
\end{frame}

\begin{frame}{Conclusion}
\begin{itemize}
\item Narrowing to the qmd-window sampling improves effective kept
samples (\(\sim24\%\rightarrow\sim45\%\)) --- geometry-acceptance result,
survives the contamination audit.
\item Clean 2026-04-20 retrain: val \(|\vec p|\) RMSE \(= 6.08\) MeV/c,
transverse gain \(p_x\) \(-51\%\), \(p_y\) \(-81\%\) vs the deleted polluted
baseline.
\item Full-event re-eval (2026-04-22) on \(5.4\)M events confirms the
\(p_x\) gain (\(-40\%\) to \(-54\%\) RMSE across overall/y-pol/z-pol); \(p_y\)
flat; \(p_z\) slightly worse because of OOD \(p_z\) tails in the eval set.
\item Clean production model: \texttt{formal\_B115T3deg/20260420\_184345/model/}.
Clean reco outputs: \texttt{eval\_20260422\_cleanmodel/}.
\item Follow-up: per-truth-bin residual breakdown (separate in-window vs
OOD \(p_z\)); consider widening training window to cover IMQMD tails; point
\texttt{run\_target\_momentum\_reco\_pipeline.sh} defaults at the clean model;
re-do hyperparameter sweep on the clean dataset.
\end{itemize}
\end{frame}

\end{document}
